% Part: counterfactuals
% Chapter: introduction
% Section: paradoxes-material

\documentclass[../../../include/open-logic-section]{subfiles}

\begin{document}

\olfileid{cnt}{int}{par}

\olsection{భౌతిక సోపాధికం యొక్క వైరుధ్యాభాసాలు}

ఆంగ్లంలోని ``ఒకవేళ \dots అయితే \dots'' వాక్యాలను
$!A \lif !B$తో సంకేతీకరించడాన్ని తొలుత
విమర్శించినవారిలో సి.~ఐ. లూయిస్ ఒకరు.
లూయిస్, వైట్‌హెడ్ మరియు రసెల్‌ల
\emph{ప్రిన్సిపియా మాథమేటికా}లో భౌతిక
సోపాధికాన్ని వాడిన తీరును విమర్శించారు;
వారు $\lif$ను ``పర్యవసిస్తుంది'' అని
ఉచ్చరించేవారు. $\lif$కు ``పర్యవసిస్తుంది''
అనే అర్థమైతే, అసత్య ప్రతిజ్ఞ~$p$ ఏదైనా
$p$ నుంచి~$q$ పర్యవసిస్తుందని
సూచిస్తుందని లూయిస్ సరిగానే అభ్యంతరం
తెలిపారు: $p$ అసత్యమైతే
$p \lif (p \lif q)$ సత్యం. అలాగే ఏదైనా
సత్య ప్రతిజ్ఞ~$q$, $p$ నుంచి~$q$
పర్యవసిస్తుందని సూచిస్తుంది: $q$
సత్యమైతే $q \lif (p \lif q)$ సత్యం.

అయితే \emph{తార్కిక పర్యవసానం}
ఒక సంయోజకం కాదనీ,
\tetoken{సూత్రాల}{!!{formula}s} లేదా
వాక్యాల మధ్య సంబంధమనీ తర్కవేత్తలకు
తెలుసు. కాబట్టి గందరగోళం నివారించడానికి
$\lif$ను ``పర్యవసిస్తుంది'' అని
చదవకూడదు.\footnote{గణితవేత్తలు,
  కంప్యూటర్ శాస్త్రవేత్తలు ఇప్పటికీ
  ``$\lif$''ను ``పర్యవసిస్తుంది'' అని
  చదవడం విస్తృతంగా పాటిస్తారు. అయితే
  లూయిస్ ఎత్తిచూపిన గందరగోళాలను
  నివారించడానికి తత్త్వవేత్తలు దాన్ని
  ``అయితే మాత్రమే'' అని పలకడానికి
  ప్రయత్నిస్తారు.} అలా చదవనంతవరకు
లూయిస్ ఆందోళన అసలు తలెత్తదు:
$p$ను అసత్య ఆంగ్ల వాక్యానికి సూచికగా
తీసుకున్నా, $p$ నుంచి~$q$
``పర్యవసించదు''. $p \Entails q$నా
అని నిర్ణయించాలంటే \emph{అన్ని}
\tetoken{సత్యమూల్య నియామకాలను}{!!{valuation}s}
పరిశీలించాలి. యాదృచ్ఛికంగా అసత్యమైన
వాక్యాన్ని $p$తో సూచించినా
$p \Entails/ q$.

అయినా లూయిస్ ఉదాహరణలాంటివి,
``ఒకవేళ \dots అయితే \dots'' అనే
వాక్యాలు కొంత విచిత్రంగా అనిపిస్తాయి:
\begin{quote}
చంద్రుడు ఆకుపచ్చ జున్నుతో తయారై ఉంటే,
$2+2=4$.
\end{quote}
ఈ నిగమనాలూ అలాగే:
\begin{quote}
  చంద్రుడు ఆకుపచ్చ జున్నుతో తయారు కాలేదు.
  కాబట్టి, చంద్రుడు ఆకుపచ్చ జున్నుతో
  తయారై ఉంటే, $2+2=4$.

  $2+2 = 4$. కాబట్టి, చంద్రుడు
  ఆకుపచ్చ జున్నుతో తయారై ఉంటే,
  $2+2=4$.
\end{quote}
అయితే ``ఒకవేళ \dots అయితే \dots''
కేవలం $\lif$ మాత్రమే అయితే, ఆ వాక్యం
నిస్సందేహంగా సత్యమవుతుంది; ఆ నిగమనాలు
నిస్సందేహంగా చెల్లుబాటవుతాయి.

మరొక ఆందోళనకు ఉదాహరణ
$(!A \lif !B) \lor (!B \lif
!A)$ అనే సర్వసత్యం. వార్తాపత్రిక
నుంచి యాదృచ్ఛికంగా రెండు సూచనాత్మక
వాక్యాలు~$S$, $T$ తీసుకుంటే,
``ఒకవేళ $S$ అయితే $T$, లేదా
ఒకవేళ $T$ అయితే~$S$'' అనే
వాక్యం సత్యంగా ఉండాలని ఇది సూచిస్తుంది.

\end{document}
